%=============================================================================!
% 16.2  WHAT DIFFRACTION SEES                                                 !
%=============================================================================!
\section{What diffraction sees}
\label{sec:ob-sq}

A street lamp seen through a fine net curtain shows a pattern of bright
spots: light passing between the threads is scattered from each gap, and
in some directions the waves from all the gaps arrive in step and add,
while in others they arrive out of step and cancel. The pattern is set by
the spacing of the threads, and from it the spacing can be read without
looking at the curtain. X-rays and neutrons scattered by the atoms of a
liquid or a crystal do the same at the scale of ångströms, and what they
measure is the structure factor, the counterpart of $g(r)$ in waves.

\subsection*{Waves scattered by atoms}

A plane wave of wave vector $\vb{k}$ has the phase $\vb{k}\cdot\vb{r}$ at
the point $\vb{r}$, its crests lying where that is a multiple of $2\pi$.
Scattered by an atom at $\vb{r}_j$ towards a distant detector, it leaves
as a wave of the same wavelength, its wave vector $\vb{k}'$ of the same
length as $\vb{k}$ and pointing at the detector. The incoming wave
reaches the atom with the phase $\vb{k}\cdot\vb{r}_j$; the path on to the
detector is shorter than the path from the origin by the length of
$\vb{r}_j$ along $\vb{k}'$, which takes off the phase
$\vb{k}'\cdot\vb{r}_j$. So the wave arrives with the phase
$\vb{k}\cdot\vb{r}_j - \vb{k}'\cdot\vb{r}_j$ plus a part that is the same
for every atom, and the waves from different atoms differ in phase by
$-\vb{G}\cdot\vb{r}_j$, with $\vb{G} = \vb{k}' - \vb{k}$ the scattering
vector. The wave from atom $j$ can be drawn as an arrow of unit length
turned through that angle, and the total as the sum of $N$ such arrows.
Turning every arrow the other way changes no length, so the angles may
be taken as $+\vb{G}\cdot\vb{r}_j$, and the total has the components
\begin{equation*}
   \mathcal{C}(\vb{G}) = \sum_j\cos(\vb{G}\cdot\vb{r}_j) , \qquad
   \mathcal{S}(\vb{G}) = \sum_j\sin(\vb{G}\cdot\vb{r}_j) ,
\end{equation*}
the sums of \cref{sec:md-ewald} with every charge 1. A detector
measures the intensity, which grows as the square of a wave's height,
here the square of the total's length; per atom it is the structure
factor
\begin{equation}
   S(\vb{G}) = \frac1N\Bigl[\mathcal{C}(\vb{G})^2 + \mathcal{S}(\vb{G})^2
   \Bigr] .
   \label{eq:ob-sq}
\end{equation}
Expanding the squares and using $\cos A\cos B + \sin A\sin B = \cos(A -
B)$ (\cref{eq:mo-addition}) writes it as a double sum,
\begin{equation*}
   S(\vb{G}) = \frac1N\sum_i\sum_j\cos\bigl(\vb{G}\cdot(\vb{r}_j -
   \vb{r}_i)\bigr) = 1 + \frac1N\sum_i\sum_{j\neq i}\cos(\vb{G}\cdot
   \vb{r}_{ij}) ,
\end{equation*}
the terms with $j = i$ each being 1. Atoms placed at random give cross
terms that average to zero, so $S(\vb{G}) = 1$; atoms of a crystal at a
vector $\vb{G}$ that makes every $\vb{G}\cdot\vb{r}_{ij}$ a multiple of
$2\pi$ give $S(\vb{G}) = N$, a Bragg peak. Such vectors form the
crystal's reciprocal lattice, built as in \cref{sec:md-ewald} on the
crystal's own repeating unit instead of the simulation cell. A liquid has
$S(\vb{G})$ above 1 near the repeat of its shells of neighbours and below
1 elsewhere.

In a periodic cell only waves that repeat with the cell are allowed, the
vectors $\vb{G} = m_1\vb{a}^* + m_2\vb{b}^* + m_3\vb{c}^*$ of
\cref{sec:md-ewald}, the shortest of length $2\pi/L$ for a cube of side
$L$. The function \texttt{structure\_factor} of \texttt{mdlab} evaluates
\cref{eq:ob-sq} on each of them, averages over the frames and over the
vectors of each length $G$, which in a liquid differ only by noise, and
returns $S(G)$ against $G$.

\subsection*{The structure factor from the arrangement}

The double sum can also be written with $g(r)$. Around each atom of a
liquid, which has no preferred direction, \cref{eq:ob-rdf} spreads its
$\rho g(r)\,4\pi r^2\,\dd r$ partners evenly over the shell, so it puts
$\rho g(r)\,\dd^3r$ in each small volume $\dd^3r$ at the separation
$\vb{r}$, and the mean of the double sum is
\begin{equation*}
   S(\vb{G}) = 1 + \rho\int g(r)\cos(\vb{G}\cdot\vb{r})\,\dd^3r
   = 1 + \rho\int\bigl[g(r) - 1\bigr]\cos(\vb{G}\cdot\vb{r})\,\dd^3r ,
\end{equation*}
the second form subtracting $\rho\int\cos(\vb{G}\cdot\vb{r})\,\dd^3r$,
which is zero over the cell for every allowed $\vb{G} \neq \vb{0}$
(\cref{sec:md-ewald}), and which makes the integrand vanish far away
where $g(r)$ is 1. On a shell of radius $r$, with $\theta$ the angle
between $\vb{r}$ and $\vb{G}$, the part at angles between $\theta$ and
$\theta + \dd\theta$ is a band of radius $r\sin\theta$ and width
$r\,\dd\theta$, of area $2\pi r^2\sin\theta\,\dd\theta$, out of $4\pi
r^2$, so the average of the cosine over the shell is, with $u =
\cos\theta$ and $\dd u = -\sin\theta\,\dd\theta$, the minus sign turning
the limits $u = 1$ to $-1$ round,
\begin{equation*}
   \frac12\int_0^\pi\cos(Gr\cos\theta)\sin\theta\,\dd\theta
   = \frac12\int_{-1}^{1}\cos(Gru)\,\dd u
   = \frac12\Bigl[\frac{\sin Gru}{Gr}\Bigr]_{-1}^{1}
   = \frac{\sin Gr}{Gr} .
\end{equation*}
The volume integral is then one over shells,
\begin{equation}
   S(G) = 1 + 4\pi\rho\int_0^\infty r^2\bigl[g(r) - 1\bigr]
   \frac{\sin Gr}{Gr}\,\dd r ,
   \label{eq:ob-sq-rdf}
\end{equation}
which depends only on the length $G$, as it must in a liquid that has no
preferred direction.

\Cref{fig:ob-sq}(a) compares the two routes for liquid argon at
\SI{135}{\kelvin}. From the positions, on 250 frames of the
\SI{2}{\nano\second} run, $S(G)$ has its first peak of 2.271 at $G =
\SI{1.985}{\per\angstrom}$. That peak lies near the repeat of the shells
of neighbours: the first three peaks of $g(r)$ lie at 3.68, 7.04 and
\SI{10.30}{\angstrom}, \SI{3.31}{\angstrom} apart on average, and
$2\pi/3.31 = \SI{1.90}{\per\angstrom}$, 4\% below the peak. Through
\cref{eq:ob-sq-rdf}, with $g(r)$ known only to half the cell, the peak is
2.126, 0.15 below, and at the smallest allowed $G$,
\SI{0.270}{\per\angstrom}, the positions give 0.058 and the transform
0.074, with ripples on either side that take it as low as $-0.05$, which
no intensity can be. Stopping the integral at $L/2$ cuts off whatever
$g(r) - 1$ holds beyond. Near the peak the oscillations of $g(r)$ beyond
the cut keep step with $\sin Gr$, so what is cut off adds up: the $g(r)$
of the 2048-atom runs, stopped at \SI{11.6}{\angstrom} and at
\SI{23.3}{\angstrom}, gives 2.116 and 2.240 there. At small $G$ the loss
is smaller but large beside $S(G)$ itself. The direct route has no such
error and is the one to trust wherever the two differ.

\subsection*{The long-wave limit and the compressibility}

As $G$ shrinks towards 0, $\sin(Gr)/Gr$ tends to 1 and
\cref{eq:ob-sq-rdf} tends to $S(0) = 1 + \rho\int[g(r) - 1]\,\dd^3r$; at
$\vb{G} = \vb{0}$ itself \cref{eq:ob-sq} gives $N$, the beam that passes
straight through, so $S(0)$ means this limit. The same integral measures
how much the number of atoms $N_{\mathrm r}$ in a region of volume
$V_{\mathrm r}$ fluctuates. The
ordered pairs of distinct atoms both inside the region number $N_{\mathrm r}(N_{\mathrm r} -
1)$. Each atom inside has on average $\rho g(r)\,\dd^3r$ partners in each
small volume at the separation $\vb{r}$, and for a region large against
the reach of $g(r) - 1$ nearly every atom is far from its surface, so
\begin{equation*}
   \ens{N_{\mathrm r}(N_{\mathrm r} - 1)}
   = \ens{N_{\mathrm r}}\,\rho\int_{V_{\mathrm r}}g(r)\,\dd^3r
   = \ens{N_{\mathrm r}}\Bigl[\rho\int\bigl[g(r) - 1\bigr]\dd^3r
   + \ens{N_{\mathrm r}}\Bigr] ,
\end{equation*}
since $\rho\int_{V_{\mathrm r}}\dd^3r$ is $\rho V_{\mathrm r} =
\ens{N_{\mathrm r}}$. Subtracting $\ens{N_{\mathrm r}}^2 -
\ens{N_{\mathrm r}}$ from both sides leaves $\operatorname{var}(N_{\mathrm r}) = \ens{N_{\mathrm r}}[1
+ \rho\int[g(r) - 1]\,\dd^3r]$, so $S(0) = \operatorname{var}(N_{\mathrm r})/
\ens{N_{\mathrm r}}$. That fluctuation follows from Chapter~14. A fixed number of
atoms held at fixed pressure has $\operatorname{var}(V) =
\kB T\ens{V}\kappa_T$ (\cref{eq:pr-volume}), and so has a group of atoms
inside a large liquid, the liquid around it acting as the barostat. Such
a group of $N_{\mathrm r}$ atoms has the density $N_{\mathrm r}/V$, which
changes by $\delta\rho/\rho = -\delta V/V$ as its volume $V$ fluctuates;
the region of fixed volume $V_{\mathrm r}$ has the density $N_{\mathrm
r}/V_{\mathrm r}$, which changes by $\delta\rho/\rho = \delta N_{\mathrm
r}/N_{\mathrm r}$ as its number fluctuates. Both are the fluctuation of the
liquid's density over the same volume, so
\begin{equation*}
   \frac{\operatorname{var}(N_{\mathrm r})}{\ens{N_{\mathrm r}}^2} =
   \frac{\operatorname{var}(V)}{\ens{V}^2} = \frac{\kB T\kappa_T}{\ens{V}} ,
   \qquad
   S(0) = \frac{\operatorname{var}(N_{\mathrm r})}{\ens{N_{\mathrm r}}} = \rho\kB T\kappa_T ,
\end{equation*}
the second following on multiplying the first by $\ens{N_{\mathrm r}} =
\rho\ens{V}$. With Chapter~14's $\kappa_T = (1.46 \pm
0.11)$\,\si{\per\giga\pascal}, from the slope $\dd P/\dd V$ of runs at fixed
volume near \SI{0.1}{\giga\pascal}, $\rho\kB T\kappa_T = 0.0554 \pm
0.0042$.

The 256-atom cell has only four lengths of wave vector below
\SI{0.6}{\per\angstrom}, too few to follow $S(G)$ towards $G = 0$, so
\cref{fig:ob-sq}(b) uses Chapter~15's runs of 1372 and 2048
atoms, two of each, whose cells allow $G$ down to
\SI{0.135}{\per\angstrom}. Each value carries the standard error of its
frames' values. A straight line in $G^2$ through the 29 values below
\SI{0.6}{\per\angstrom}, fitted as in \cref{sec:cv-size} with each squared
miss divided by the square of that value's error, reaches $S(0) = 0.0589$
at $G = 0$. The sum of its squared misses in units of the errors is 72,
where scatter alone gives $27 \pm 7$ for 29 values less the two fitted
numbers, so the errors understate the misses by about $\sqrt{72/27} =
1.63$, and the line's error, 0.0002, enlarged by that factor becomes
0.0004. The two values, $0.0589 \pm 0.0004$ from the positions and
$0.0554 \pm 0.0042$ from the change of pressure with volume, lie 0.8
combined errors apart: the long waves of the structure and the liquid's
stiffness against squeezing measure the same compressibility. The test
is loose in one respect: Chapter~14 took its slope at a volume 2.2\%
smaller than this run's, where the denser liquid is the stiffer.

The same integral fixes the far value of $g(r)$ in a fixed cell
(\cref{sec:ob-rdf}). With the divisor $N\rho$, an atom's partners over
the whole cell number $N - 1$, so $\rho\int_V[g(r) - 1]\,\dd^3r = -1$. In a
cell large enough for $g(r) - 1$ to die away well inside it, the part of
the integral near the atom is $S(0) - 1$, and the remaining $-S(0)$ is
spread over the rest of the cell, where $g(r) - 1$ is then $-S(0)/\rho V =
-S(0)/N$\cite{Lebowitz1961}: $1 - 1/N$ for atoms at random, whose $S(0)$
is 1, and $1 - 0.059/N$ for this liquid. In the corners of the 2048-atom
cell the partners fall short of $\rho$ by a fraction $(3.3 \pm
0.5)\times10^{-5}$, which is $S(0)/N$ for $S(0) = 0.067 \pm 0.009$, the
value above within its error.

\begin{figure}[htb]
\centering
\includegraphics{ch16_observables/sq}
\caption{(a) The structure factor of liquid argon at \SI{135}{\kelvin}
from the positions, averaged over the wave vectors of the 256-atom cell
of each length up to $12\times2\pi/L = \SI{3.24}{\per\angstrom}$ (points),
against the
transform of $g(r)$ cut at half the cell (dashed). (b) $S(G)$ at the
smallest wave vectors of the 1372- and 2048-atom runs of Chapter~15, two
runs of each, with their standard errors, and a straight line in $G^2$
fitted below \SI{0.6}{\per\angstrom} (dashed), against $\rho\kB T\kappa_T$
with Chapter~14's compressibility (band, its error).}
\label{fig:ob-sq}
\end{figure}

\begin{keyidea}
The structure factor $S(G) = [\mathcal{C}^2 + \mathcal{S}^2]/N$ is the
intensity, per atom, of waves scattered with the change of wave vector
$\vb{G}$, and what diffraction measures. It is the transform of
$g(r) - 1$ by $\sin(Gr)/Gr$, which a cut at half the cell spoils, by
0.15 at the peak and by a large fraction at small $G$, so it is computed
from the positions on the cell's own wave vectors. Its limit at $G = 0$ is $\rho\kB T\kappa_T$.
\end{keyidea}

\notebook{16}{compare S(G) from the positions and from g(r)}

\begin{selfcheck}
For a crystal at $T = 0$, what is $S(\vb{G})$ at a vector $\vb{G}$ of
its reciprocal lattice, and at a vector of the cell that is not one of
them?
\end{selfcheck}
